\section{Annual evaluation and state transfer}
Each annual case follows this procedure:
\begin{enumerate}
\item Build the continuous UTC price timeline and assign Europe/London date and profile-slot labels. Load the quarter-hour workload and tranche arrays.
\item Initialise UPS/TES energy and the four thermal states once, at the beginning of the year; the initial unfinished-work ledger is empty.
\item For local day $d$, select its 92, 96 or 100 core intervals and append twelve further intervals. Load the carried state and unfinished cohorts, and generate new arrivals throughout the horizon.
\item Solve the reference or coordinated model, and record feasibility, objective bounds and accounting checks. Both core and look-ahead contain the full facility demand.
\item Commit only core controls and processing. Accumulate signed-price settlement for those intervals, save their ending physical state and compute unfinished work from committed processing only.
\item Pass the ending state and unfinished cohorts to the next local day. Reoptimise all future controls. Repeat through 31 December.
\end{enumerate}
If $\boldsymbol x$ contains the two energy states and four temperatures, physical transfer is $\boldsymbol x_{d+1}^{\mathrm{start}}=\boldsymbol x_d^{\mathrm{core,end}}$. There is no daily storage equality restoring the initial state. Storage and workload transfer are separate: an energy state alone cannot reconstruct which jobs remain eligible for execution.

For annual committed set $\mathcal Y$, reported settlement is
\begin{equation}
C=\sum_{t\in\mathcal Y}\pi_tP_t^{\mathrm{grid}}\Delta t/1000.
\end{equation}
This counts every 2025 interval once. The final continuation contains ordinary January demand as well as outstanding December work, so its entire cost is not added to 2025 expenditure. The central ledger contains 0.787115 CPU-hours at midnight after 31 December, and the retained continuation serves that work within its permitted execution windows. The final continuation's feasibility does not imply that its full dispatch is committed as an additional annual day.

The implementation records each day's controls, physical states, cohort residuals, solver status and audits. Central operation also records actual processing by cohort and interval, permitting event initialisation from the exact unfinished-work ledger. Recorded daily objectives include look-ahead decisions, whereas reported annual costs include committed decisions only. A daily MILP gap is therefore not an optimality certificate for a single full-year optimisation problem.

\section{Representative-day and flexibility-event procedures}
\subsection{Day selection}
Daily features from central coordinated operation are: mean price, price standard deviation, minimum and maximum price, negative-price share, mean and peak grid import, mean total CPU utilisation, daily settlement cost, and opening UPS, TES and cold-aisle states. Each feature is centred on its mean across all annual daily observations and scaled by its population standard deviation; a zero standard deviation is replaced by one. For candidate day $d$, the score is
\begin{equation}
s_d=\left[\frac{1}{12}\sum_{j=1}^{12}
\left(\frac{v_{dj}-\bar v_j}{\sigma_j}\right)^2\right]^{1/2}.
\end{equation}
Candidates have 96 intervals and local date before 29 December 2025. The smallest score is selected, with date used to break a tie. This gives 22 September, score 0.23481769. The complete candidate table accompanies the package. The rule identifies a central multivariate illustration; it does not weight days by event availability, or guarantee a typical duration envelope.

\subsection{Opening state and recovery}
For each hourly $t_0$, load the undisturbed central trajectory and exact remaining work at that boundary. Pre-event execution is fixed by the stored history; the event solve starts at $t_0$. New arrivals from that point are generated with their original class rules. Any existing cohort retains only its remaining work and remaining eligible intervals.

For duration $\tau$, the trial ends at $H=t_0+\tau+3$ h. Impose a grid deviation of $\Delta P$ within $\pm0.1$ kW during the event, with no grid-target constraint during the recovery period. At $H$, require UPS/TES energy no smaller than the undisturbed states, each temperature within 0.05$^{\circ}$C of its undisturbed value, and each cohort residual no larger than the undisturbed residual plus $10^{-7}$ CPU-hours. Workload eligibility, storage modes and all physical limits hold throughout. Surplus stored energy is permitted at recovery; exact equality of storage states is not required.

\subsection{Search, certification and dispatch interpretation}
There are 24 start times, nine negative requests ($-100$ to $-500$ kW in 50 kW steps) and three positive requests (25, 50, 75 kW). For each pair, the assessment cap in intervals is the smaller of 48 and the number of intervals remaining in the local day. The procedure first tests that cap. If unsuccessful, it searches logarithmically for a feasible duration and checks the following interval. Trials in this assessment have a 60 s solver limit. The eight unresolved adjacent trials are retained as lower-bound results.

A trial with an audited feasible incumbent is evidence that its duration can be delivered under the model. A trial without an incumbent can be either proven infeasible or unresolved; these outcomes are retained separately. During the search, an unsuccessful trial can narrow the region explored even when it is unresolved. Consequently the results should be interpreted as the feasible durations found by this procedure, rather than as an exhaustive exploration of every duration.

The output flag \path{boundary_verified} covers both a feasible assessment cap and a proven-infeasible adjacent candidate. It is not, by itself, a proof that all longer durations are infeasible. Since extending an event changes the recovery boundary and undisturbed terminal reference, monotonic feasibility with duration has not been established for this formulation. The main paper therefore reports certified feasible durations. It retains plus signs for the eight unresolved adjacent cases and explains the cap convention separately. Unmarked cells do not acquire a stronger global-maximality interpretation merely because the next tested interval was infeasible. Blank cells mean that this search certified no positive duration.

The accepted dispatch minimises signed event-plus-recovery electricity cost subject to the request. Where a solve terminates with a nonzero objective gap, it is an accepted incumbent for that objective, not a proven minimum-cost allocation. Even a proven optimum need not be unique. Thus component panels illustrate feasible coordinated mechanisms and exact electrical accounting; they are neither a unique attribution nor an independent valuation of the assets. The dashed target ends at the certified event duration. Grey space aligns panels of different lengths and contains no plotted recovery dispatch; recovery is checked against the terminal constraints. A plus sign beside a panel's duration denotes an unresolved adjacent trial, consistently with the heatmap.
